% arara: xelatex: { shell: yes, synctex: yes }
\RequirePackage[l2tabu, orthodox]{nag}
\documentclass{article}

% -- Texto e codificação
\usepackage[utf8]{inputenc}
\usepackage[T1]{fontenc}
\usepackage[portuguese]{babel}
\usepackage{microtype}
\usepackage{xcolor}
\usepackage{anyfontsize}
\usepackage{pdflscape}

% -- Tipo de letra
\usepackage{newpxtext}
\usepackage{eulervm}
\usepackage{nimbusmono}

%\usepackage{lmodern}
%\usepackage[T1]{fontenc}

% -- Definir margens do documento
\usepackage{geometry}
\geometry{nomarginpar,
    tmargin = 2.5cm,
    bmargin = 2.5cm,
    lmargin = 2.5cm,
    rmargin = 2.5cm}
%\usepackage{showframe}

% -- Cabeçalho e rodapé
\usepackage{fancyhdr}
\fancyhf{}
\fancyhead[L]{\textsc{}}
\fancyhead[R]{\textsc{}}
\fancyfoot[C]{\thepage}
\pagestyle{fancy}
\renewcommand\headrulewidth{0.2pt}
\renewcommand\footrulewidth{0pt}

% -- Funções matemáticas extra
\usepackage{mathtools}
\usepackage{siunitx}

% -- Símbolos extra
\usepackage{amssymb}
\usepackage{textcomp}
\usepackage{gensymb}
\usepackage{cancel}

% -- Referências
\usepackage{hyperref}
\hypersetup{colorlinks,
	linkcolor	= {red!50!black},
	citecolor	= {green!30!black},
	urlcolor	= {blue!80!black}}

% -- Bibliografia
\usepackage[
	backend = biber,
	style = alphabetic,
	sorting = ynt
	]{biblatex}
\usepackage{fvextra}
\usepackage{csquotes}

% --  Definições de imagens
\usepackage{graphicx}
\graphicspath{{graphics/}}
\usepackage{caption}
\usepackage{subcaption}
\usepackage{afterpage}
\usepackage{tabularx}

% -- Desenhar circuitos elétricos e lógicos
\usepackage{tikz}
\usepackage{pgfplots}
\usetikzlibrary{arrows.meta,positioning}
\pgfplotsset{compat=1.5}
\pgfplotsset{table/search path = {data}}
\pgfplotsset{	/pgf/number format/use comma,}

% -- Integrar código fonte
\usepackage{minted}

\usepackage{circuitikz}
\usetikzlibrary{positioning}

\newcommand{\matr}[1]{\mathbf{#1}}

\addbibresource{bibliography.bib}

\begin{document}

\begin{titlepage}

\begin{center}
	\vspace*{0.1\textheight}
	\includegraphics[width = 0.4\linewidth, trim = {172.4pt 202.7pt 172.6pt 201.4pt}, clip]{IST_A_CMYK_POS}
	
	\vspace*{0.1\textheight}
%	\vspace{\fill}
	{\huge\bfseries Real-Time Cooperative Decentralized Control of a Smart Office Illumination System}
	
	\vspace*{0.03\textheight}
	{\Large Distributed Real-Time Control Systems Final Project}
	
	\vspace*{43.5mm}
	{\Large \begin{tabular}{l r} Daniel de Schiffart & \texttt{81479} \\ João Gonçalves & \texttt{81040} \\ Francisco Castro & \texttt{78655}\end{tabular}}
	
	\vspace{\fill}
	{\large Distributed Real-Time Control Systems}
	
	\vspace*{0.01\textheight}
	{\Large Master's Degree in Aerospace Engineering}
	
	\vspace*{0.01\textheight}
	{\large Instituto Superior Técnico}
\end{center}

\end{titlepage}
\setcounter{page}{1}

{\hypersetup{linkcolor = black} \tableofcontents}

\section{Introdução}

\begin{tikzpicture}
	\node at (0,0) (ard1) [anchor=south west,draw,minimum height=3,inner sep=1] {Arduino 1};
	\node (ard2) [left=of ard1,anchor=south west,draw,minimum height=3,inner sep=1] {Arduino 2};
\end{tikzpicture}

\section{Raspbian Startup}
\begin{enumerate}
	\item Open console
	\item \texttt{sudo ifconfig enp2s0f0 10.0.0.2 netmask 255.255.255.0}
	\item Configure by hand on Ubuntu wired connection settings
	\begin{enumerate}
		\item Details -> Disable \textit{Connect Automatically}
		\item IPv4 -> \textit{IPv4 Method} to \textit{Manual}
		\item Set address to \texttt{10.0.0.2} and gateway to \texttt{255.255.255.0}
	\end{enumerate}
	\item Run \texttt{ssh pi\@10.0.0.1} and use \texttt{raspberry} as password
\end{enumerate}

\textbf{Importante:} antes de desligar o Raspberry Pi, correr \texttt{sudo shutdown -h now} para evitar corrupção do sistema operativo.

\section{Optimização do Sistema Distribuído}

\begin{gather*}
	u^* = \textrm{argmin}\,J(x, \mu) \\
	l_i = \sum^n_{j=1} d_j k_{ij} + o_i \geq L_i, \qquad 0 \leq d_i \leq 100, \quad \forall i \\
	\matr{l} = \matr{K}\matr{d} + \matr{o} \geq \matr{L}, \qquad 0 \leq d \leq 100
\end{gather*}
Se as condições aqui listadas não forem cumpridas, o problema é declarado \textit{unfeasible}. A função de custo é
\begin{gather*}
	f(d) = c_1d_1 + c_2d_2 + \ldots + c_nd_n = \matr{C}^T\matr{d}
\end{gather*}
onde $\matr{c}$ é o vetor dos custos de energia.

Ver os cantos do domínio admissível e saber o ponto com menor custo. Assim, obtém-se $(d_1,d_2)$. Mais à frente adicionamos maximização de longevidade de lâmpadas.
\begin{figure}[ht]
	\centering
	\begin{tikzpicture}
		\coordinate (axis) at (0,0);
		\node at (0,5) [label={180:$d_2$}] {};
		\node at (5,0) [label={-90:$d_1$}] {};
		\draw [->] (axis) ++(0,-0.1) -- (0,5);
		\draw [->] (axis) ++(-0.1,0)-- (5,0);
		\draw [dashed,color={green!50!black}] (axis) (0,4) -- (4,4) -- (4,0);
	\end{tikzpicture}
\end{figure}

\begin{tikzpicture}
	\tikzstyle{arduino} = [draw,minimum height=80,inner sep=30];
	\tikzstyle{wire} = [];
	\node at (0,0) (ard1) [arduino] {\Large Arduino 1};
	\node (ard2) [right=4cm of ard1,arduino] {\Large Arduino 2};
	\node (rpi) [above=2cm of ard2,arduino] {\Large Raspberry Pi};
	\path (ard1.north east) to node (ard1_scl) [auto, swap, align=right] {\ttfamily SCL} ++(0,-3);
	\path (ard1.north east) to node (ard1_sda) [auto, swap, align=right] {\ttfamily SDA} ++(0,-4);
	\path (ard1.north east) to node (ard1_gnd) [auto, swap, align=right] {\ttfamily GND} ++(0,-5);
	\path (ard2.north west) to node (ard2_scl) [auto, align=left] {\ttfamily SCL} ++(0,-3);
	\path (ard2.north west) to node (ard2_sda) [auto, align=left] {\ttfamily SDA} ++(0,-4);
	\path (ard2.north west) to node (ard2_gnd) [auto, align=left] {\ttfamily GND} ++(0,-5);
	\path (rpi.north west) to node (rpi_19) [auto, align=left] {\ttfamily 19} ++(0,-3);
	\path (rpi.north west) to node (rpi_18) [auto, align=left] {\ttfamily 18} ++(0,-4);
	\path (rpi.north west) to node (rpi_gnd) [auto, align=left] {\ttfamily GND} ++(0,-5);
	\draw (ard1_scl) -- (ard2_scl) (ard1_sda) -- (ard2_sda) (ard1_gnd) -- (ard2_gnd);
	\path (ard1_scl) ++(1.7,0) coordinate (coord1);
	\path (ard1_sda) ++(2.4,0) coordinate (coord2);
	\path (ard1_gnd) ++(3.1,0) coordinate (coord3);
	\draw (coord1) to [short,*-,R] (coord1 |- rpi_19) to [short,-] (rpi_19.west);
	\draw (coord2) to [short,*-,R] (coord2 |- rpi_18) to [short,-] (rpi_18.west);
	\draw (coord3) to [short,*-] (coord3 |- rpi_gnd) to [short,-] (rpi_gnd.west);
\end{tikzpicture}

\pagebreak
\printbibliography

\end{document}
